
%======================================================================
\section{v9: what is new, and the response to the priority list}
\label{sec:whatv9}
%======================================================================

v9 is built mechanically from the read-only v8 source
(\loc{code/build_v9.py}). Every v8 statement is kept unless this
section or Table~\ref{tab:open} says it has changed. v9 has new
evidence class marks: \emph{Executed in v9} means a script under
\loc{code/v9/}, a log under \loc{logs/v9/}, and a SHA-256 in
Appendix~\ref{app:v9artifacts}. The frozen trees XVI-P, XVI-Q and
XVI-R were hashed before and after the v9 runs. Both digests equal
the v8 values (Section~\ref{sec:v9hash}). Files v1--v8 were not
edited.

\begingroup\small
\begin{longtable}{@{}L{.05\linewidth}L{.27\linewidth}L{.20\linewidth}L{.40\linewidth}@{}}
\caption{v9 verdicts on the nine priority items.}\label{tab:v9verdict}\\
\toprule
\# & Item & Verdict & Evidence / obstruction\\
\midrule
\endfirsthead
\toprule
\# & Item & Verdict & Evidence\\
\midrule
\endhead
1 & Hidden mixed D01 terms & computed; v6 \emph{restored} &
Proposition~\ref{prop:v9TT}: pure TT, vanishes on the lattice.
\(\Delta J(\alpha)\) at 14 angles (Table~\ref{tab:v9dJ}).
IR log: Proposition~\ref{prop:v9log}.\\
2 & \(r^2\)-weighted mixed sum, D04 part of 2nd \(\phi\phi\) &
\OPEN, precise obstruction & Section~\ref{sec:v9D04}: the jet densities
are not exported, and same-\(k\) pairing is not IR-summable. No
closed form is claimed.\\
3 & CK/CB dictionary; \(e_5,e_6\); closed forms &
dictionary computed; \(e_5,e_6\) \emph{excluded with proof} at constant profile &
Theorem~\ref{thm:v9ghost}: all 6 constant-profile Ward functionals
hold; \(\ker\pi_KB\) is one-dimensional.\\
4 & 20 jet-level Ward combinations, full projection &
\OPEN, precise obstruction & Section~\ref{sec:v9jet}.\\
5 & 18-dim engine map beyond 6 & \OPEN, precise obstruction &
Section~\ref{sec:v9jet}.\\
6 & Evanescent vs frozen \(\tilde g\) & \emph{excluded}, precise obstruction &
Proposition~\ref{prop:v9evan}.\\
7 & Vector-type keys & log parts computed; finite parts \OPEN &
Proposition~\ref{prop:v9vector}.\\
8 & Generic-\(\alpha\) rank certificate &
computed on chambers; exceptional set \PARTIAL &
Theorem~\ref{thm:v9generic}.\\
9 & New sheets & read; one newer sheet & Section~\ref{sec:v9sheets};
\loc{new_oct18g/} written.\\
\bottomrule
\end{longtable}
\endgroup

\paragraph{Changes to v8 results.}
(a) The v8 status ``\CONDITIONAL{} on the hidden D01 term'' for the v6
\(hh\) branch difference is lifted. The value is unchanged and is
now proved on the lattice (Proposition~\ref{prop:v9TT}).
(b) The four CK/CB-supported Ward combinations, which v8
identified only conditionally, are now checked exactly on the
catalogue. This uses two fitted, \(\alpha\)-independent ghost
normalizations (Theorem~\ref{thm:v9ghost}).
(c) The v8 open item ``\(e_5\) and \(e_6\)'' is resolved negatively:
the constant-profile data determine only \(e_6-e_2\) and
\(e_5+e_6/2\) beyond the v8 combinations.
(d) The incomplete-beta reduction used in an internal v8/v9 draft
module for the F-branch \(k=-1\) mode (\loc{I_F}) contained an
exponent error (\(x^{1-\mu}\) for \(x^{\mu-1}\)). It was not used
in any printed v8 number. It was removed, and every v9 value uses
the closed-form quadrature \loc{I_gen}.
No other v8 number changes.

%======================================================================
\section{v9: the hidden D01 pairs}
\label{sec:v9hidden}
%======================================================================

\subsection{Closed form of the mixed trace}

Let \(A_\mu\) be the radial Bessel operator of order \(\mu\ge0\) on
the cone channel, and let
\[
 g(x;\mu,\nu)=\Tr\bigl[e^{-(1-x)A_\mu}\,r^{-2}\,e^{-xA_\nu}\bigr],
 \qquad 0<x<1,
\]
with \(d=\mu-\nu\) and \(h=(\mu+\nu)/2\).

\begin{proposition}[Mixed trace; derived and executed]
\label{prop:v9g}
\[
 g(x;\mu,\nu)=\frac{\Gamma(h)\Gamma(h+1)}{4\,\Gamma(\mu+1)\Gamma(\nu+1)}\;
 {}_2F_1\!\bigl(\tfrac d2+1,\tfrac d2;\mu+1;x\bigr)\;
 {}_2F_1\!\bigl(1-\tfrac d2,-\tfrac d2;\nu+1;1-x\bigr).
\]
In particular \(g(x;\mu,\mu)=1/(4\mu)\), and
\(\int_0^1g\,\dd x=1/(2(\mu+\nu))\).
\end{proposition}

\begin{proof}[Derivation]
Expand both Bessel kernels in their power series and integrate the
radial variables term by term. This gives
\(x^\mu(1-x)^\nu\) times
\(F_4(h,h+1;\mu+1,\nu+1;x^2,(1-x)^2)\). Since
\(h+(h+1)-(\mu+1)+1=\nu+1\), Bailey's reduction
\(F_4(a,b;c,a+b-c+1;X(1-Y),Y(1-X))={}_2F_1(a,b;c;X)\,{}_2F_1(a,b;a+b-c+1;Y)\)
applies with \(X=x\) and \(Y=1-x\). Euler's transformation of each
factor removes the prefactor \(x^\mu(1-x)^\nu\) and gives the
stated parameters. The integral identity is an independent check
of the v8 Duhamel lemma, Lemma~\ref{lem:v8duhamel}.
\loc{code/v9/g_formula_check.py} confirms both the closed form and
the integral against direct quadrature of the kernels. The
residual, \(4.7\times10^{-3}\), is the truncation tail of the
quadrature at \(R=60\), not a discrepancy.
\end{proof}

\subsection{The hidden pairs are pure TT}

With \(\hat\rho=6x^2-10x+3\) and
\(\hat R=x(1-x)(3-2x)\), we have \(\hat R'=\hat\rho\). For a mode
pair \((\mu,\nu)\), integration by parts gives
\[
 I(\mu,\nu)=\int_0^1\hat\rho\,f
 =-(\mu^2-\nu^2)\int_0^1\hat R\,g .
\]
\loc{code/v9/modefun.py} evaluates this in two independent ways
for \(\nu=\mu\pm2\): an exact Gamma series summed by
Euler--Maclaurin, and direct quadrature of
Proposition~\ref{prop:v9g}. The two agree to 30 digits.

\begin{proposition}[TT structure; proved]
\label{prop:v9TT}
On the D01 keys \((4,4),(4,7),(7,4),(7,7),(5,5)\), the hidden
\((0,2)\)/\((2,0)\) pairs contribute
\[
 -\frac{\kappa^2P^2J}{8\pi}\Bigl[(h_{11}-h_{22})(h'_{11}-h'_{22})
 +4h_{12}h'_{12}\Bigr],
\]
with \(J=\sum_{\rm modes}I\) in the stated mode list. On the v6 lattice
(\(h_{11}=h_{22}=f\), no \(h_{12}\)) the bracket is identically zero.
Hence the v6 \(hh\) branch difference
\(-\nu\kappa^2P^2f_0f_0'/(4\pi)\) holds without the v8 condition.
\end{proposition}

\begin{proof}
The key-by-key coefficients come from
\loc{logs/v8/slice/hidden_mixed.json} and the catalogue
normalization
\((P^2\kappa^2/16\pi)/(\ii(q_0^2-q_3^2)/8\pi)\). Key \((4,4)\)
receives \(-\kappa^2P^2J/(8\pi)\); keys \((4,7)\) and \((7,4)\) carry
the opposite sign, and key \((5,5)\) carries four times key
\((4,4)\). Collecting terms gives the bracket. The lattice
projection of \loc{invariant_data} uses
\(\{0{:}0,3{:}1,9{:}2,4{:}3,7{:}3\}\). It sets \(h_{11}=h_{22}\)
and drops index~5, so it maps the bracket to zero for every value
of \(J\).
\end{proof}

\subsection{Infrared structure and the finite branch difference}

\begin{proposition}[IR logarithm; executed asymptotics]
\label{prop:v9log}
As \(\mu\to\infty\),
\(I(\mu,\mu+2)=\tfrac13+\tfrac1{30\mu}-\tfrac1{10\mu^2}+\cdots\) and
\(I(\mu,\mu-2)=-\tfrac13+\tfrac1{30\mu}+\tfrac1{10\mu^2}+\cdots\)
(coefficients up to \(\mu^{-8}\) are in \loc{logs/v9/asym02.json}).
Under same-\(k\) pairing, the constant terms cancel between \(k\)
and \(-k\), but the \(1/\mu\) terms add to \(\alpha/(15k)\).
Therefore \(J\) diverges like \((\alpha/15)\log\Lambda\). In
general, a pair \((s,t)\) contributes
\(\tfrac{\alpha}{2}(s-t)^2\,x(x-1)\log\Lambda\) to \(C_{s,t}(x)\).
This is the frame/flux mismatch of a long-range \(r^{-2}L\)
perturbation.
\end{proposition}

So the absolute TT values are not pure apex coefficients. Their
finite part depends on the scheme (for example \(k\)-\(\zeta\)).
The CR and F branches differ only in the \(S=2\), \(k=-1\) mode,
whose order is \(-\nu_0\) (CR) or \(+\nu_0\) (F), with
\(\nu_0=2-1/\alpha\). The difference
\[
 \Delta J(\alpha)=I(1/\alpha,\,1/\alpha-2)-I(1/\alpha,\,2-1/\alpha)
\]
is therefore finite and does not depend on the regularization.
The corresponding catalogue difference is
\(-\kappa^2P^2\Delta J/(8\pi)\) times the TT bracket, and it is
invisible on the lattice.

\begin{center}
\scriptsize
\begin{longtable}{@{}lllll@{}}
\caption{The finite TT branch difference (executed, \loc{code/v9/hidden02_run.py}; CR mode by series and by quadrature, agreement \(<10^{-30}\)).}\label{tab:v9dJ}\\
\toprule
\(\alpha\) & \(\nu_0\) & \(I_{\rm CR}\) (\(k=-1\)) & \(I_{\rm F}\) (\(k=-1\)) & \(\Delta J\)\\
\midrule
@@DJTABLE@@
\bottomrule
\end{longtable}
\end{center}
At \(\alpha=1/2\), \(\nu_0=0\) and \(\Delta J=0\) exactly. The sign
of \(\Delta J\) is the sign of \(-\nu_0\) at every computed angle.
As \(\alpha\to1\), the CR mode degenerates (\(\mu+\nu\to0\)); the
values decrease toward zero (\(\Delta J(99/100)\approx-4.9\times10^{-3}\)).
PSLQ was not attempted. We claim no elementary
closed form, since each value is an integral of a product of two
\({}_2F_1\) factors.

%======================================================================
\section{v9: the CB/CK dictionary and the six constant-profile Ward functionals}
\label{sec:v9ghost}
%======================================================================

\subsection{Dictionaries}

We use the same method as the \(hh\) map (Section~\ref{sec:v8map}).
Unit \(m=0\) atoms \(L=(1,0),(0,1)\) for \(c,\bar c\), and
\(LL\)/\(PM\) atoms for \(K\), are evaluated by \loc{eval_sym}
(reflection only, so \(\Gamma_{\rm sym}=2\Gamma\)) on the complete
\(6\times6\) momentum grid. A dictionary is accepted only if the
catalogue bilinear (D06 for \(\langle\bar c\,c\rangle\), D13 for
\(\langle K c\rangle\)) lies \emph{exactly} in the span of the
evaluated \(K\)-candidates on all columns
(\loc{code/v9/cbck_map.py}, \loc{code/v9/ck_free.py}).

\begin{proposition}[CB and CK dictionaries; executed, exact]
\label{prop:v9dict}
With \(q_0=\pm\ii q_\tau\), \(q_3=q_z\):
\begin{enumerate}[label=(\roman*)]
\item CB: exactly the four sign choices with opposite \(\ii\)-signs
for \(\bar c_0\) and \(c_0\) are accepted. All four give
\(y_{\rm CB}=\ii(-R/8,\,0)\,\kappa^2/\pi\), with \(R=3z_0+4z_1\).
\item CK: with the \(K\) components left free, the solution is
unique: \(K^{00}=-LL_{\tau\tau}\), \(K^{03}=\mp2\ii LL_{\tau z}\), and
\(y_{\rm CK}=\ii\bigl((1-\alpha)/16,\,R/8,\,0\bigr)\kappa^2/\pi\).
The \(PM\) scale is undetermined because the catalogue's PM keys
are computed zeros.
\end{enumerate}
Both hold at \((3/10,\CR)\), \((3/10,\Fr)\), \((7071/10^4,\CR)\) and
\((7071/10^4,\Fr)\). Here
\(R=(17\alpha-7)(\alpha-1)/(12\alpha)\).
\end{proposition}

\subsection{The six functionals}

The annihilator of \(\pi_KZ\) is printed in
\loc{apex_cert_*_vectors.json} and is identical at all four
points. In the order
\(K=(0,1,2,5,6,12,18,19,26,27,29,32)\) it consists of
\(W_1=y_0+y_1\), \(W_2=y_2\), \(W_3=-2y_0+y_{19}\),
\(W_4=-2\ii y_0+4\ii y_5+4\ii y_6+y_{26}\),
\(W_5=-2\ii y_{18}+y_{27}\), and \(W_6=2\ii y_0+y_{29}\). (v8
quoted only the real parts of these vectors.)

\begin{theorem}[Ghost-sector Ward check; executed, exact]
\label{thm:v9ghost}
Put \(\lambda_{\rm CB}=-4\ii\) and \(\lambda_{\rm CK}=-8\), fixed
once and independent of \(\alpha\) and of the branch. Rescale the
CB and CK coordinates by these factors. Then all six \(W_i\)
vanish on the catalogue at all four certified points. Solving
\(\sum_ie_i\pi_KB_i=y_K\) on all twelve coordinates gives
\[
 e_3=0,\quad e_1=\tfrac{1-\alpha}4,\quad
 e_6-e_2=-\tfrac R2,\quad
 e_5+\tfrac{e_6}2=\frac{\alpha^2-1}{48\alpha},
\]
\[
 e_4-\tfrac{e_6}2=\frac{119\alpha^2-120\alpha+25}{192\alpha}\ (\CR),\qquad
 e_4-\tfrac{e_6}2=\frac{7(17\alpha-7)(\alpha-1)}{192\alpha}=\tfrac{7R}{16}\ (\Fr).
\]
The kernel of \(\pi_KB\) is exactly
\(\Span\{(0,1,0,\tfrac12,-\tfrac12,1)\}\). Therefore \(e_5\) and
\(e_6\) are not determined separately by constant-profile data.
\end{theorem}

\begin{proof}
\loc{code/v9/e_solve.py} carries out the linear algebra exactly; its output is
\loc{logs/v9/e_solve.json}. The checks have different strengths.
\(W_3\) and \(W_6\) hold because both of their entries vanish.
\(W_5\) defines \(\lambda_{\rm CB}\), since both of its terms are
proportional to \(R\).
\(W_4\) is a genuine check. Its \(hh\) part is
\(4\ii(b_6+b_7)=\ii(c_0-c_1)=\ii(1-\alpha)/2\), coming from
D01--D03 and D11--D12. Its CK part is
\(\lambda_{\rm CK}\,\ii(1-\alpha)/16\), coming from D13. The two
have the same \(\alpha\)-shape, so one constant
\(\lambda_{\rm CK}=-8\) works at every \(\alpha\) and on both
branches. The \(e\)-values follow from the twelve equations. The
closed forms are the rational identities displayed above,
checked at the four points.
\end{proof}

\noindent\emph{Scope.} The normalizations \(\lambda\) are fitted,
not derived from the catalogue's BRST conventions. The closed
forms hold on the certified points. For general generic \(\alpha\)
they follow from the closed forms of \(z_0,z_1,z_2\) together with
the observed \(\alpha\)-independence of \(\pi_KB\) and of the
annihilator, which was verified at four points but not proved.
Status: \PARTIAL{} for general \(\alpha\), exact on the four points.
The v8 identity \(e_1=(1-\alpha)/4\) is now obtained twice, from
the \(hh\) block and from the CK coordinate~26.

%======================================================================
\section{v9: chambers and a generic-\texorpdfstring{\(\alpha\)}{alpha} rank certificate}
\label{sec:v9chamber}
%======================================================================

The structural key of a jet \((C,n,m)\) is
\((C,n,\epsilon,j,\operatorname{sgn}(n/\alpha+s))\), where
\(m=\epsilon\nu+2j\). Two angles lie in the same \emph{chamber} when
every gh0 (\(T=4\)) and gh\(-1\) (\(T=3\)) sector has the same set
of structural candidates. In a chamber, every ST row entry is a
fixed rational function of \(\alpha\).

\begin{proposition}[Walls; executed and proved]
\label{prop:v9walls}
In \((1/4,1)\) the walls are exactly \(\{1/3,1/2,2/3\}\), on both
branches. CR has the same structural signature on all four
chambers. F has one signature on \((1/4,1/3)\cup(1/3,1/2)\) and a
different one, with \(N=38\) at the certified points, on
\((1/2,2/3)\cup(2/3,1)\).
\end{proposition}

\begin{proof}
A non-generic candidate requires \(\pm2n/\alpha\in\mathbb Z\) for
some admissible pair \((n,-n)\), or a sign or \(\epsilon\) change.
The jet budget gives \(\nu\le6\) and \(n\le4\). So a wall has the
form \(\alpha=2n/M\) with \(M\le12\), or \(\alpha=n/|s|\). Every
such rational has denominator at most 12.
\loc{code/v9/chamber_scan.py} evaluates the structural signature
at every rational with denominator at most 40 in \((1/4,1)\),
which is exhaustive by this bound. It finds only the three walls
listed (\loc{logs/v9/chamber_cache2.json}).
\end{proof}

\begin{theorem}[Generic rank over \(\mathbb Q(\alpha)\)]
\label{thm:v9generic}
On every CR chamber in \((1/4,1)\), and on the F chambers in
\((1/4,1/2)\), the ST rank over \(\mathbb Q(\alpha)\) is exactly
\(26\), with \(\dim Z=8\) generically. On the F chambers in
\((1/2,1)\), the generic rank \(r\) satisfies \(29\le r\le30\).
\end{theorem}

\begin{proof}
\emph{Lower bound.} The rank is lower semicontinuous under
specialization within a chamber. So it suffices to have one point
per chamber with a nonzero minor of the required size. Exact
\(\mathbb Q(\ii)\) minors exist at \(3/10\) and \(7071/10^4\)
(Theorem~\ref{thm:v8cert}). A nonzero minor mod \(p\) also implies
a nonzero rational minor, and the modular runs at \(41/100\),
\(309/500\), \(4/5\), \(8501/10^4\) and at the scan points of
Table~\ref{tab:v9scan} provide such minors.
\emph{Upper bound.} The identity \(\sigma_0\sigma_{-1}=0\) holds
identically, so \(B\subset Z\). We have \(s\phi=0\), so \(PP\subset Z\).
\(B\) has no PP components, and \(\dim B\ge6\) generically, by lower
semicontinuity from the certified points. Hence \(\dim Z\ge8\),
which gives \(r\le N-8\): \(26\) for \(N=34\) and \(30\) for
\(N=38\). The ninth cocycle seen at the F points with \(N=38\) is
not shown to persist generically, so \(29\le r\le30\) there.
\end{proof}

\begin{center}
\scriptsize
\begin{longtable}{@{}lll@{}}
\caption{Modular scan (\loc{code/v9/apex_scan.py}, an unchanged copy of the historical driver with a new output directory; prime \(10^9+9\), randomized complete-rank sampling). Every rational \(p/q\in(1/4,1)\) with \(q\le20\).}\label{tab:v9scan}\\
\toprule
branch & result & angles\\
\midrule
@@SCANTABLE@@
\bottomrule
\end{longtable}
\end{center}

\noindent\emph{Exceptional set.} Inside a chamber, the rank drops
exactly on the common zero set of the maximal minors. That set is
finite, but the minor polynomial was not computed: it would need a
symbolic-\(\alpha\) elimination of unknown degree. So the
exceptional set is \(E=\{1/3,1/2,2/3\}\cup E_{\rm minor}\). The
scan shows that \(E_{\rm minor}\) contains no rational with
denominator at most 20, but \(E_{\rm minor}\) itself is \OPEN.

%======================================================================
\section{v9: precise obstructions}
\label{sec:v9obs}
%======================================================================

\subsection{The D04 part of the second \texorpdfstring{\(\phi\phi\)}{phi phi} value}
\label{sec:v9D04}

The D04 internal lines are a graviton, with channel weights
\(\mathrm{pw}\), and \(\phi\) (\(S=0\)). At constant profile only
\(z_0\) contributes, and the same-channel part is the
\(x\)-independent \(D\)-moment
\(\Tr[r^2e^{-A}]\). The second \(\phi\phi\) value needs the
\(r^2\)-weighted mixed traces
\(\Tr[e^{-(1-x)A_\mu}r^2e^{-xA_\nu}]\) with transverse-momentum
numerators and the \(H_I\) intertwining. There are two
obstructions:
\begin{enumerate}[label=(\alph*)]
\item the zero-transverse-momentum slice does not export these
jet-level densities, so the integrand is not available as data;
\item the frame/flux mismatch of Proposition~\ref{prop:v9log} carries
over with an extra \(r^2\). Same-\(k\) pairing is then not
IR-summable. The mismatch gives the log of the \(r^{-2}\) case
there, and power counting gives a power divergence with the
\(r^2\) weight (\emph{argued, not executed}).
\end{enumerate}
So a rational reconstruction in \(\alpha\) has nothing
well-defined to reconstruct until a scheme for the mixed
same-\(k\) sum is declared. Status: \OPEN; no closed form is
claimed, and no checks at \(\alpha=1\) or \(1/2\) are claimed.

\subsection{Jet-level Ward combinations and the engine map}
\label{sec:v9jet}

Twenty of the 26 Ward combinations, and twelve of the eighteen
\(hh\) engine coordinates, involve \(m>0\) jets (\(r^2\) profiles)
or transverse \(P\)/\(M\) components. The dictionary method of
Sections~\ref{sec:v8map} and \ref{sec:v9ghost} would need the
catalogue's values on \(m=2\) and weight-\(\pm1\) atoms. The
declared catalogue has no such values: it is the zero-transverse,
constant-profile slice. They would come from the
\(r^2\)/\(D_b\)-weighted traces, which are obstructed as in (b)
above. The full 34-dimensional projection is therefore available
only on the twelve \(K\) coordinates, where all six functionals now
hold (Theorem~\ref{thm:v9ghost}). That projection has rank 6, and
its angle and branch dependence enters only through \(z_0,z_1,z_2\).
Status: \OPEN.

\subsection{Evanescent coefficients versus the frozen \texorpdfstring{\(\tilde g\)}{g-tilde} rule}

\begin{proposition}[Excluded, with obstruction; executed]
\label{prop:v9evan}
The frozen engine (\loc{XVI-P/code/p5_entry/dred.py},
\loc{eta_trace_Gt}, read-only) admits exactly one evanescent
pattern, \(c\sum_r\eta^{rr}\tilde G_{rr}\to c(4-d)\), and declares
every other pattern open. Of the 285 v8 slice-DRED keys
(\loc{logs/v8/slice/evanescent_combined.json}), 164 contain
\(\tilde G\) monomials of degree 2 or 3. All of them contain only
the longitudinal \(\tilde G_{00},\tilde G_{03},\tilde G_{33}\), and
only 31 have a linear part of partial-trace form
\(c(\tilde G_{00}-\tilde G_{33})\)
(\loc{code/v9/evan_match.py}). The v8 coefficients therefore lie
outside the class the frozen rule can evaluate. A match would need
the transverse \(\tilde G_{11},\tilde G_{22}\) components, which
the slice does not export, and a rule for non-trace monomials.
\end{proposition}

\subsection{Vector-type keys}

\begin{proposition}[IR logs of the hidden keys; executed]
\label{prop:v9vector}
Let \(L_{\rm key}=\sum_{(s,t)}(s-t)^2\int_0^1\rho_{(s,t)}x(x-1)\)
over D01 and D03 (\loc{code/v9/vector_logs.py}), and put
\(c=\ii P^2/(30\pi)\). Then the vector keys \((1,6),(2,8),(6,1),(8,2)\)
have \(L=0\). The keys \((1,1)\) and \((2,2)\) have \(L=+c\), and
\((6,6)\) and \((8,8)\) have \(L=-c\). The TT keys give
\(-c,+c,+c,-c,-4c\) on \((4,4),(4,7),(7,4),(7,7),(5,5)\).
\end{proposition}

The vector-key log is
\(c\,P^2\sum_{a=1,2}\eta^{\mu\nu}h_{\mu a}h'_{\nu a}\) up to key
normalization. It does not cancel within D01\(+\)D03, so the
vector keys inherit the scheme dependence of
Proposition~\ref{prop:v9log}. Whether profile-dependent Ward
partners remove it cannot be decided from the slice. The finite
parts of the vector keys stay \OPEN{} and are excluded, not zero.

%======================================================================
\section{v9: frozen-tree digests and the new registry sheets}
\label{sec:v9hash}
\label{sec:v9sheets}
%======================================================================

The tree digests of XVI-P, XVI-Q and XVI-R, in the shell
serialization of Section~\ref{sec:v8hash}, were recomputed before
the first v9 run and after the last one:
@@HASHROWS@@
They are equal to the v8 values. The v8 source and PDF were
hashed before and after as well
(\loc{logs/v9/v8_hashes_before.sha}, \loc{logs/v9/v8_hashes_after.sha}).

\paragraph{Registry.} The registry was fetched again on 18 October
2026. It lists 258 identifiers, one more than the 257 in
\loc{new_oct18f/}: the new identifier is \loc{I-C-Ex2}. It was
stored under \loc{new_oct18g/} with SHA-256 sums. The six sheets
of \loc{new_oct18f/} and \loc{I-C-Ex2} were read as cited inputs
(abstracts and all theorem statements; \loc{logs/v9/sheets_extract.txt}):
\begin{itemize}
\item \emph{I-B-G-S-B} \cite{CheungIBGSBv9} is the only sheet that
bears on a v9 input. For a declared class of nonnegative tensor
connection forms, it proves Mosco and strong-resolvent selection of
the componentwise Friedrichs realization, and states that for
neutral spin~2 and \(\alpha>1/2\) the limit excludes the
curvature-regular branch. This bears on the CR/F branch input: it
favours F \emph{within that declared core class}. It does not
derive the graviton operator, and the sheet itself stresses that a
smooth metric and its holonomy do not fix the perturbation
operator. So both branches remain inputs here, and the CR results
are kept as conditional on CR.
\item \emph{I-B-ECHO-v2} \cite{CheungEchoV2} gives the scalar
apex-diffracted Green operator and finite report instruments. It
is not an apex-residue input; compare Example~\ref{ex:image}.
\item \emph{X-J-ST-2} \cite{CheungXJST2} gives the free photon
reference (Maxwell--Friedrichs with a Friedrichs scalar ghost) and
charged scalar coefficients. It is a free, source-test reference
and imports no interacting graviton Ward identity.
\item \emph{Hawk05-ST, I-C-v2, I-C-Ex-v2, I-C-Ex2}
\cite{CheungHawk05ST,CheungICv2,CheungICExv2,CheungICEx2} are
finite access and recovery bounds (\(1-1/d\), \(1-k/d\)). They
belong to the instrument side of the boundary of
Section~\ref{sec:readout} and are not used for O.1.
\end{itemize}
